\metatitle{Distance and Tutte polynomials}
\metadescription{A distance polynomial for colored matroids, its tropical
minor recurrence, its joint-rank expression, and the classical Tutte polynomial.}

\section[distanceAndTutte]{Distance and Tutte polynomials}

The color-indexed minors suggest looking for an analogue of the Tutte
polynomial. Simply weighting the words of a colored matroid is too coarse:
if the color quotas are $(r_1,\dotsc,r_k)$, then
\[
 \sum_{u\in\mathcal F}\prod_{e\in E}x_{u_e}
 =|\mathcal F|x_1^{r_1}\dotsm x_k^{r_k}.
\]
To retain the geometry of the family, we instead measure the distance from
every coloring of $E$ to the allowed words.

\subsection[distancePolynomial]{The distance polynomial}

Let $\mathcal F$ be a colored matroid on $E$, with $n=|E|$ and colors $[k]$. For
$w\in[k]^E$, define its \defin{distance to $\mathcal F$} by
\[
 \defin{d_{\mathcal F}(w)}
 \coloneqq
 \min_{u\in\mathcal F}|\{e\in E:w_e\ne u_e\}|.
\]

Write $\mathbf x=(x_1,\dotsc,x_k)$.

\begin{definition}[Distance polynomial]
The \defin{distance polynomial of $\mathcal F$} is
\[
 \defin{D_{\mathcal F}(\mathbf x;t)}
 \coloneqq
 \sum_{w\in[k]^E}
 t^{d_{\mathcal F}(w)}\prod_{e\in E}x_{w_e}.
\]
The variables $x_1,\dotsc,x_k$ record the colors of $w$, while $t$ records
the number of entries that must be changed to reach $\mathcal F$.
\end{definition}

\subsection[rankDistance]{Distance from joint rank}

The \href{rank-flats.htm#jointRank}{joint rank} counts the largest number
of prescribed assignments that a word of $\mathcal F$ can realize.
Thus distance measures the assignments left unmatched.

\begin{proposition}[Rank and distance]
For every word $w\in[k]^E$, put $W_c=B(w,c)$. Then
\[
 d_{\mathcal F}(w)=n-R(W_1,\dotsc,W_k).
\]
Consequently, the distance polynomial is
\[
 \begin{gathered}
 D_{\mathcal F}(\mathbf x;t)\\
 =\sum_{w\in[k]^E}t^{\,n-R(W_1,\dotsc,W_k)}
 \prod_{c=1}^k x_c^{|W_c|}.
 \end{gathered}
\]
\end{proposition}

\begin{proof}
For each $u$, its number of disagreements with $w$ is $n$ minus its
number of agreements. The latter is $\sum_c|B(u,c)\cap W_c|$.
Minimizing disagreements therefore subtracts the maximum agreement
count $R(W_1,\dotsc,W_k)$ from $n$.
\end{proof}

The minor operation gives a recurrence for the distance itself.

\begin{proposition}[Tropical minor recurrence]
Fix $e\in E$, let $w_e=a$, and write $w'$ for the restriction of $w$ to
$E\setminus\{e\}$. Then
\[
 d_{\mathcal F}(w)
 =
 \min_{\substack{c\in[k]\\\mathcal F/(e,c)\ne\varnothing}}
 \left(
  \mathbf 1_{a\ne c}+d_{\mathcal F/(e,c)}(w')
 \right),
\]
where $\mathbf 1_{a\ne c}$ is $0$ if $a=c$ and $1$ otherwise.
\end{proposition}

\begin{proof}
Partition the words $u\in\mathcal F$ according to the color $c=u_e$. For
such a word, its disagreement with $w$ at $e$ contributes
$\mathbf 1_{a\ne c}$, and the remaining disagreements are measured by the
minor $\mathcal F/(e,c)$. Taking the minimum over the nonempty branches gives
the formula.
\end{proof}

This is a min-plus, or tropical, recurrence. The proof uses only Hamming
distance, so the same recurrence holds for every nonempty family of words.

\subsection[basicProperties]{Basic properties}

\begin{proposition}
Let $\mathcal F$ have quotas $(r_1,\dotsc,r_k)$ on an $n$-element set.
Then the distance polynomial has the following properties.
\begin{enumerate}
\item If $\mathcal F=\mathcal F_1\oplus\mathcal F_2$, then
\[
 D_{\mathcal F}(\mathbf x;t)
 =D_{\mathcal F_1}(\mathbf x;t)
  D_{\mathcal F_2}(\mathbf x;t).
\]
\item Its two immediate specializations are
\[
 D_{\mathcal F}(x_1,\dotsc,x_k;0)
 =|\mathcal F|x_1^{r_1}\dotsm x_k^{r_k}
\]
and
\[
 D_{\mathcal F}(x_1,\dotsc,x_k;1)
 =(x_1+\dotsb+x_k)^n.
\]
\item After setting every $x_c=1$, the coefficient of $t^j$ counts the
colorings at distance $j$ from $\mathcal F$. In particular, the constant
coefficient is $|\mathcal F|$, and the degree is
$\max_{w\in[k]^E}d_{\mathcal F}(w)$.
\item If every word in $\mathcal F$ has color $c$ at $e$, then
\[
 \begin{gathered}
 D_{\mathcal F}(\mathbf x;t)\\
 =\left(x_c+t\sum_{a\ne c}x_a\right)
  D_{\mathcal F/(e,c)}(\mathbf x;t).
 \end{gathered}
\]
\end{enumerate}
\end{proposition}

\begin{proof}
Distances add across a direct sum, which proves the first identity. At $t=0$
only the words of $\mathcal F$ remain, while at $t=1$ all $k^n$ colorings
remain. Setting the color variables to $1$ gives the distance distribution.
For the last identity, changing the color at $e$ costs one unless the new
color is $c$, independently of the remaining positions.
\end{proof}

Simultaneously permuting the colors of $\mathcal F$ and the corresponding
variables leaves the polynomial unchanged. For two colors, exchanging the
colors corresponds to matroid duality.

\subsection[twoColors]{Two colors and the Tutte polynomial}

Suppose that $M$ is a rank-$r$ matroid on an $n$-element set and that
$\mathcal F$ is its family of basis words: color $1$ marks a basis and color
$0$ its complement.

\begin{proposition}
For a binary word $w$, let $A$ be its set of color-$1$ positions. Then
\[
 d_{\mathcal F}(w)
 =\bigl(r-r_M(A)\bigr)+\bigl(|A|-r_M(A)\bigr).
\]
Consequently,
\[
 \begin{gathered}
 D_{\mathcal F}(x_0,x_1;t)\\
 =x_1^r x_0^{n-r}
 T_M\left(1+t\frac{x_0}{x_1},
           1+t\frac{x_1}{x_0}\right).
 \end{gathered}
\]
Thus the two-color distance polynomial is a homogeneous form of the classical
Tutte polynomial.
\end{proposition}

\begin{proof}
For a basis $B$, the Hamming distance from $w$ to its basis word is
$|A\mathbin{\triangle}B|$. The largest possible value of $|A\cap B|$ over all
bases $B$ is $r_M(A)$, and hence
\[
 \min_{B\in\mathcal B(M)}|A\mathbin{\triangle}B|
 =|A|+r-2r_M(A).
\]
This is the sum of the corank $r-r_M(A)$ and the nullity
$|A|-r_M(A)$. Substituting this expression in the definition of
$D_{\mathcal F}$ and comparing with
\[
 \begin{gathered}
 T_M(X,Y)=\\
 \sum_{A\subseteq E}
  (X-1)^{r-r_M(A)}(Y-1)^{|A|-r_M(A)},
 \end{gathered}
\]
gives the stated identity. This is the corank--nullity expansion of the Tutte
polynomial; see Oxley \cite{Oxley2011}.
\end{proof}

\subsection[linearRecurrence]{Why the recurrence is tropical}

The classical Tutte polynomial has a linear deletion--contraction recurrence.
The analogous identity
\[
 D_{\mathcal F}(\mathbf x;t)
 \stackrel{?}{=}
 \sum_{\substack{c\in[k]\\\mathcal F/(e,c)\ne\varnothing}}
 x_cD_{\mathcal F/(e,c)}(\mathbf x;t)
\]
does not hold for every colored matroid.

\begin{example}
Take
\[
 \mathcal F=\{\mathtt{123},\mathtt{132},\mathtt{213}\}
\]
and remove position $e=2$. Each pair of words differs along a single color
cycle, so this is a colored matroid. Its three minors at $e$ are
\[
 \begin{gathered}
 \mathcal F/(2,1)=\{\mathtt{23}\},\\
 \mathcal F/(2,2)=\{\mathtt{13}\},\\
 \mathcal F/(2,3)=\{\mathtt{12}\}.
 \end{gathered}
\]
The coefficient of $x_2x_3^2t$ in $D_{\mathcal F}$ is $3$, contributed by
$\mathtt{233}$, $\mathtt{323}$, and $\mathtt{332}$. The coefficient on the
right-hand side of the proposed recurrence is $2$.

The word $\mathtt{112}$ shows what goes wrong pointwise. Its distance from
$\mathcal F$ is $1$, attained by $\mathtt{132}$. If we insist on keeping its
color $1$ at position $2$, the only available word is $\mathtt{213}$, and
after removing position $2$ the distance from $\mathtt{12}$ to
$\mathtt{23}$ is $2$. A nearest word need not lie in the minor selected by
the color of $w_e$.
\end{example}

\begin{question}
For which colored matroids, or which positions $e$, does the linear identity
above hold?
\end{question}
